\section{Introdução}\label{sec:introduction}

O projeto de estruturas de madeira exige representar um material variável por um conjunto limitado de propriedades. O uso de classes de resistência oferece uma forma operacional de realizar essa representação. Entretanto, a diferença entre propriedades de uma população ensaiada e os valores atribuídos à sua classe pode afetar de maneiras distintas a previsão de deslocamentos, as verificações resistentes e as dimensões escolhidas. Quantificar essas consequências permite avaliar quando a caracterização experimental acrescenta informação relevante ao projeto.

\textcite{Wolenski2020} estudaram 40 espécies nativas do grupo das folhosas e avaliaram as relações $f_{c0,k}=0{,}77f_{t0,k}$ e $E_{c0}=E_{t0}$ previstas na NBR 7190:1997. Os autores compararam médias por análise de variância e ajustaram modelos de regressão entre propriedades de compressão e tração paralelas às fibras. Sua interpretação foi favorável às relações normativas, com melhor ajuste dos modelos geométricos entre as alternativas avaliadas. O conjunto publicado inclui médias, desvios-padrão, coeficientes de variação e resistências características, permitindo examinar diferenças entre espécies além das relações médias globais.

A aplicação a pontes rurais tem antecedentes de campo: \textcite{Criado2024} inspecionaram dez pontes com pranchas serradas e vigas roliças e discutiram patologias e pré-dimensionamento. Esse tipo de levantamento evidencia a necessidade de relacionar hipóteses de cálculo às condições dos elementos empregados, embora não corresponda a uma avaliação controlada da substituição de propriedades experimentais por valores de classe.

A otimização de pontes também possui antecedentes que vão além da apresentação de ferramentas computacionais. \textcite{Marti2013} combinaram decisões geométricas, materiais e armaduras no projeto de pontes pré-moldadas de concreto protendido, utilizando busca heurística para estudar alternativas de menor custo. Em pontes de madeira laminada colada, \textcite{SimoesNegrao2005} formularam um procedimento de projeto ótimo baseado em confiabilidade para passarelas estaiadas, integrando análise estrutural e múltiplos critérios. Esses trabalhos sustentam o uso da otimização como instrumento de comparação, mas seus sistemas e objetivos diferem da investigação da representação por classe em pontes de madeira roliça.

O presente estudo parte dessa caracterização para formular uma pergunta estrutural: quanto a substituição das propriedades experimentais pelas propriedades de classe altera o desempenho e o consumo de material de uma ponte? As relações entre compressão e tração investigadas na fonte não respondem diretamente a essa pergunta. A ausência de diferença significativa entre médias de grupos tampouco demonstra, por si só, equivalência dentro de uma margem relevante para o projeto ou limita os desvios de cada espécie em relação à sua classe.

Pontes de madeira para estradas vicinais constituem a aplicação considerada. Nessa tipologia, a rigidez participa da verificação de deslocamento, enquanto as resistências participam das verificações últimas. O vão, as ações e as restrições geométricas podem alterar a importância relativa dessas verificações. Essa influência deve ser avaliada, sem assumir antecipadamente que a flecha governe ou que ocorra uma transição de estado limite ao longo da faixa estudada.

Três operações precisam ser distinguidas. A primeira é a obtenção da resistência característica, que pode alterar a classe atribuída. A segunda é a substituição da rigidez experimental pelo valor da classe. A terceira é o redimensionamento, que modifica a geometria e pode compensar ou ampliar diferenças observadas na avaliação de uma geometria fixa. Misturar essas operações impede identificar o mecanismo responsável pela diferença entre projetos.

O objetivo é quantificar esses efeitos a partir da base de \textcite{Wolenski2020}. Os objetivos específicos são: (i) auditar as propriedades e a classificação reportadas na fonte; (ii) comparar o enquadramento com resistência característica publicada e com uma estimativa simplificada a partir da média; (iii) medir o efeito da representação da rigidez sobre deslocamentos de uma mesma estrutura; (iv) comparar o consumo de material após redimensionamento sob critérios comuns; e (v) verificar a influência do vão, das restrições ativas e das hipóteses necessárias às propriedades não ensaiadas na fonte. A otimização é empregada como instrumento de comparação, e não como contribuição central.

A contribuição pretendida é a ligação rastreável entre dado experimental, regra de classificação e consequência estrutural, mediante comparações pareadas. Ao manter as demais hipóteses comuns e reavaliar a geometria do projeto por classe, o protocolo permite distinguir diferença de consumo de material e atendimento ao limite de serviço. Sua generalidade reside no procedimento de comparação; os resultados quantitativos permanecerão condicionados à população, ao sistema estrutural e às ações avaliados.
